% Source: https://artofproblemsolving.com/community/c3h1656531_2018_ms_mte_sprint_all
% Stable identity: aops-topic-1656531-q16
The exterior angles of a convex polygon add up to $360^\circ$, so a regular $k$-gon has exterior angle $\dfrac{360^\circ}{k}$ and interior angle $D(k)=180^\circ-\dfrac{360^\circ}{k}$ (for the hexagon: $120^\circ$ and $60^\circ$).

\begin{asy}
usepackage("amsmath");
size(4cm);
defaultpen(fontsize(10pt));
pair[] P;for(int k=0;k<6;++k)P.push(dir(60*k));
draw(P[0]--P[1]--P[2]--P[3]--P[4]--P[5]--cycle,linewidth(0.9));
draw(P[0]--P[0]+0.75*(P[0]-P[5]),dashed);
draw(arc(P[0],0.28,degrees(P[5]-P[0]),degrees(P[1]-P[0])),red);
label("$120^\circ$",P[0]+0.48*dir(180),red+fontsize(8pt));
draw(arc(P[0],0.38,degrees(P[1]-P[0]),degrees(P[0]-P[5])),blue);
label("$60^\circ$",P[0]+0.55*dir(90),blue+fontsize(8pt));
label("\mbox{}",(min(currentpicture,user=true)+max(currentpicture,user=true))/2,black); // keeps the text after the figure black when LaTeX typesets the labels inline
\end{asy}

Hence
\[
D(m)+D(n)=360^\circ-360^\circ\left(\dfrac{1}{m}+\dfrac{1}{n}\right),
\]
and the condition becomes $\dfrac{1}{3}<\dfrac{1}{m}+\dfrac{1}{n}\le\dfrac{1}{2}$. Checking:
\begin{itemize}
\item (a) $\dfrac{1}{3}+\dfrac{1}{6}=\dfrac{1}{2}$: yes ($D(3)+D(6)=60^\circ+120^\circ=180^\circ$).
\item (b) $\dfrac{1}{5}+\dfrac{1}{8}=\dfrac{13}{40}<\dfrac{1}{3}$: no ($243^\circ$).
\item (c) $\dfrac{1}{4}+\dfrac{1}{9}=\dfrac{13}{36}$, between $\dfrac{1}{3}$ and $\dfrac{1}{2}$: yes ($230^\circ$).
\item (d) $\dfrac{1}{6}+\dfrac{1}{8}=\dfrac{7}{24}<\dfrac{1}{3}$: no ($255^\circ$).
\item (e) $\dfrac{1}{8}+\dfrac{1}{11}=\dfrac{19}{88}<\dfrac{1}{3}$: no.
\end{itemize}
The pairs are \boxed{\textbf{(a)}\ (3,6)\text{ and }\textbf{(c)}\ (4,9)}.
